\documentclass[10pt,a4paper]{amsart}
\usepackage{amsmath,amssymb,amsthm,graphicx}
% Retained preamble for the Stacks Project chapter morphisms.tex (Morphisms of Schemes).
% The chapter file pinned-original-source.tex carries only \input{preamble} on its first line, so no
% preamble text lives inside the chapter itself.  The Stacks Project uses one single shared preamble for
% every chapter; the snapshot stages it verbatim in sources/stacks-chow/preamble.tex.  The block below is
% copied mechanically, line for line, from that file's lines 14-18 (the Xy-pic setup) and lines 154-217
% (the native theorem declarations together with the mathematical macros).  Nothing was added, removed or
% redefined.  The wrapper class and the amsmath/amssymb/amsthm/graphicx base packages are supplied by the
% assembler, so the class/\externaldocument part of the shared preamble is not repeated here.
% For commutative diagrams we use Xy-pic
\usepackage[all]{xy}

% We use 2cell for 2-commutative diagrams.
\xyoption{2cell}
% Theorem environments.
%
\theoremstyle{plain}
\newtheorem{theorem}[subsection]{Theorem}
\newtheorem{proposition}[subsection]{Proposition}
\newtheorem{lemma}[subsection]{Lemma}

\theoremstyle{definition}
\newtheorem{definition}[subsection]{Definition}
\newtheorem{example}[subsection]{Example}
\newtheorem{exercise}[subsection]{Exercise}
\newtheorem{situation}[subsection]{Situation}

\theoremstyle{remark}
\newtheorem{remark}[subsection]{Remark}
\newtheorem{remarks}[subsection]{Remarks}

\numberwithin{equation}{subsection}

% Macros
%
\def\lim{\mathop{\mathrm{lim}}\nolimits}
\def\colim{\mathop{\mathrm{colim}}\nolimits}
\def\Spec{\mathop{\mathrm{Spec}}}
\def\Hom{\mathop{\mathrm{Hom}}\nolimits}
\def\Ext{\mathop{\mathrm{Ext}}\nolimits}
\def\SheafHom{\mathop{\mathcal{H}\!\mathit{om}}\nolimits}
\def\SheafExt{\mathop{\mathcal{E}\!\mathit{xt}}\nolimits}
\def\Sch{\mathit{Sch}}
\def\Mor{\mathop{\mathrm{Mor}}\nolimits}
\def\Ob{\mathop{\mathrm{Ob}}\nolimits}
\def\Sh{\mathop{\mathit{Sh}}\nolimits}
\def\NL{\mathop{N\!L}\nolimits}
\def\CH{\mathop{\mathrm{CH}}\nolimits}
\def\proetale{{pro\text{-}\acute{e}tale}}
\def\etale{{\acute{e}tale}}
\def\QCoh{\mathit{QCoh}}
\def\Ker{\mathop{\mathrm{Ker}}}
\def\Im{\mathop{\mathrm{Im}}}
\def\Coker{\mathop{\mathrm{Coker}}}
\def\Coim{\mathop{\mathrm{Coim}}}

% Boxtimes
%
\DeclareMathSymbol{\boxtimes}{\mathbin}{AMSa}{"02}

%
% Macros for moduli stacks/spaces
%
\def\QCohstack{\mathcal{QC}\!\mathit{oh}}
\def\Cohstack{\mathcal{C}\!\mathit{oh}}
\def\Spacesstack{\mathcal{S}\!\mathit{paces}}
\def\Quotfunctor{\mathrm{Quot}}
\def\Hilbfunctor{\mathrm{Hilb}}
\def\Curvesstack{\mathcal{C}\!\mathit{urves}}
\def\Polarizedstack{\mathcal{P}\!\mathit{olarized}}
\def\Complexesstack{\mathcal{C}\!\mathit{omplexes}}
% \Pic is the operator that assigns to X its picard group, usage \Pic(X)
% \Picardstack_{X/B} denotes the Picard stack of X over B
% \Picardfunctor_{X/B} denotes the Picard functor of X over B
\def\Pic{\mathop{\mathrm{Pic}}\nolimits}
\def\Picardstack{\mathcal{P}\!\mathit{ic}}
\def\Picardfunctor{\mathrm{Pic}}
\def\Deformationcategory{\mathcal{D}\!\mathit{ef}}

% --- appended 2026-10-12 (Chief): the project own definitions of the reference/slogan/history
% display environments, copied mechanically from the shared Stacks preamble lines 8-12
% (\usepackage{verbatim} plus the three \newenvironment lines, which the chapter preamble keeps
% verbatim from the project preamble). Nothing added, removed or redefined: these environments
% carry no mathematical content and the project itself typesets them as comments. Without these
% three lines a counted statement whose own text uses them fails with Environment slogan undefined.
\usepackage{verbatim}
\newenvironment{reference}{\comment}{\endcomment}
\newenvironment{slogan}{\comment}{\endcomment}
\newenvironment{history}{\comment}{\endcomment}
\usepackage{cleveref}
\crefname{originallemma}{original lemma}{original lemmas}
\Crefname{originallemma}{Original lemma}{Original lemmas}
\makeatletter
\newlabel{algebra-lemma-2-3-ring-map}{{0EUH}{source}{}{sourcepointer-1}{}}
\newlabel{algebra-lemma-2-3-ring-map@cref}{{[originallemma][0][]0EUH}{[1][0][]source}{}{}{}}
\makeatother
\crefname{originallemma}{original lemma}{original lemmas}
\Crefname{originallemma}{Original lemma}{Original lemmas}
\makeatletter
\newlabel{algebra-lemma-characterize-finite-presentation}{{00QO}{source}{}{sourcepointer-2}{}}
\newlabel{algebra-lemma-characterize-finite-presentation@cref}{{[originallemma][0][]00QO}{[1][0][]source}{}{}{}}
\makeatother
\crefname{originallemma}{original lemma}{original lemmas}
\Crefname{originallemma}{Original lemma}{Original lemmas}
\makeatletter
\newlabel{algebra-lemma-image-dense-generic-points}{{00FL}{source}{}{sourcepointer-3}{}}
\newlabel{algebra-lemma-image-dense-generic-points@cref}{{[originallemma][0][]00FL}{[1][0][]source}{}{}{}}
\makeatother
\crefname{originallemma}{original lemma}{original lemmas}
\Crefname{originallemma}{Original lemma}{Original lemmas}
\makeatletter
\newlabel{algebra-lemma-p-ring-map}{{00I5}{source}{}{sourcepointer-4}{}}
\newlabel{algebra-lemma-p-ring-map@cref}{{[originallemma][0][]00I5}{[1][0][]source}{}{}{}}
\makeatother
\crefname{originallemma}{original lemma}{original lemmas}
\Crefname{originallemma}{Original lemma}{Original lemmas}
\makeatletter
\newlabel{algebra-lemma-ring-colimit-fp}{{00QN}{source}{}{sourcepointer-5}{}}
\newlabel{algebra-lemma-ring-colimit-fp@cref}{{[originallemma][0][]00QN}{[1][0][]source}{}{}{}}
\makeatother
\crefname{originallemma}{original lemma}{original lemmas}
\Crefname{originallemma}{Original lemma}{Original lemmas}
\makeatletter
\newlabel{algebra-lemma-when-colimit}{{07C3}{source}{}{sourcepointer-6}{}}
\newlabel{algebra-lemma-when-colimit@cref}{{[originallemma][0][]07C3}{[1][0][]source}{}{}{}}
\makeatother
\crefname{originallemma}{original lemma}{original lemmas}
\Crefname{originallemma}{Original lemma}{Original lemmas}
\makeatletter
\newlabel{lemma-reduction-universal-homeomorphism}{{054M}{source}{}{sourcepointer-7}{}}
\newlabel{lemma-reduction-universal-homeomorphism@cref}{{[originallemma][0][]054M}{[1][0][]source}{}{}{}}
\makeatother
\crefname{originallemma}{original lemma}{original lemmas}
\Crefname{originallemma}{Original lemma}{Original lemmas}
\makeatletter
\newlabel{lemma-subalgebra-inherits}{{0CN7}{source}{}{sourcepointer-8}{}}
\newlabel{lemma-subalgebra-inherits@cref}{{[originallemma][0][]0CN7}{[1][0][]source}{}{}{}}
\makeatother
\crefname{originalproposition}{original proposition}{original propositions}
\Crefname{originalproposition}{Original proposition}{Original propositions}
\makeatletter
\newlabel{proposition-universal-homeomorphism}{{0CNE}{source}{}{sourcepointer-9}{}}
\newlabel{proposition-universal-homeomorphism@cref}{{[originalproposition][0][]0CNE}{[1][0][]source}{}{}{}}
\makeatother
\title{A ring map inducing a universal homeomorphism on spectra is a filtered colimit of finitely presented algebras inducing universal homeomorphisms, and the same with residue field isomorphisms}
\author{The Stacks Project Authors}
\date{stacks-project:morphisms.tex:local-snapshot-20261012; GNU FDL 1.2 or later}
\begin{document}
\maketitle
\noindent\textit{Editorial note.} All delivered mathematical text is the author's own, copied from the pinned Stacks Project chapter file pinned-original-source.tex; no statement or proof step was rewritten, reordered or removed.\par\medskip
\noindent\textit{Editorial note.} Following the multi-claim convention for this source, source\_url is the source-level URL https://stacks.math.columbia.edu/ and the source label lemma-universal-homeo-limit is carried in claim\_key and here rather than in a per-claim tag URL.\par\medskip
\noindent\textit{Editorial note.} The preamble is the chapter's shared Stacks Project preamble, retained as sources/stacks-morphisms/retained-preamble.tex: a line-for-line excerpt of the shared Stacks preamble staged with this snapshot. No macro was invented or redefined, and this card adds no preamble text of its own.\par\medskip
\noindent\textit{Editorial note.} Statement span 11910-11923 and proof span 11925-11981 come from the halfA pool card operations/pool-cards/stacks-morphisms-halfA.json and were re-read in the pinned chapter file before this card was written.\par\medskip
\noindent\textit{Editorial note.} Deferral recorded: part (2) of the statement is not proved in the counted span; the author writes that its proof is exactly the same as the proof of part (1). This is treated as a deferral rather than as a rejected candidate, because the author supplies the argument for part (1) in full and asserts the identity of the two proofs, which is materially different from declining to prove a clause (contrast the candidate lemma-relative-seminormalization, skipped today, where the author writes that the proof of part (2) is omitted with no indication of an argument). Corpus practice supports this reading: 131 accepted cards already in focused-seeds contain the phrase 'exactly the same' inside their counted proof spans.\par\medskip
\noindent\textit{Editorial note.} Tag self-check per LANE\_BRIEF 10.5 (5): the label lemma-universal-homeo-limit resolves in the project's own tag list to canonical tag 0EUJ (line 18554 of the index, full label morphisms-lemma-universal-homeo-limit). reference/public-items.jsonl holds no row whose primary source\_file is morphisms.tex, so the intersection with the same-chapter public tag ids is empty and the candidate passes.\par\medskip
\noindent\textit{Editorial note.} The gate was run as: python operations/agent-authoring/tools/precard\_gate.py --source morphisms.tex --label lemma-universal-homeo-limit --start 11910 --end 11923 --proof-start 11925 --proof-end 11981 --local-id 34090, giving RESULT=OK with all four span guards passing, no public interval row for this chapter and occupancy FREE for morphisms.tex 11910 11981. Spans were recomputed live from the label in the pinned chapter file; the located text matched the pool byte counts (statement 536 characters, proof 2314 characters after whitespace removal). The candidate lies outside the two windows of this chapter handed to lane-s1 (8000-10000 and 12000-14000).\par\medskip
\noindent\textit{Editorial note.} 9 cross-references occur inside the counted spans; they are listed in external\_references as navigation pointers only.\par\medskip
\noindent\textit{Allowed standard prerequisites (editorial).}\begin{itemize}
\item The criterion for writing an algebra as a filtered colimit of finitely presented algebras, cited from the algebra chapter. Cited at proof lines 11926-11927. Sets up the shape of the proof of (1): it suffices to factor every finitely presented algebra mapping to B through a member of the colimit.
\item A subalgebra inherits the property of inducing a universal homeomorphism on spectra, cited from the same chapter. Cited at proof lines 11931-11932. Shows that the image of a finitely presented algebra mapping to B still induces a universal homeomorphism, so the criterion can be applied to it.
\item Writing an algebra as a filtered colimit of finitely presented algebras with surjective transition maps, and factoring a map from a finitely presented algebra through a member of such a colimit, cited from the algebra chapter. Cited at proof lines 11933-11939. Produces the presentation of the image as a colimit of finitely presented algebras and an index through which the given finitely presented algebra factors.
\item Reduction of a universal homeomorphism to the associated reduced rings, and the dense image criterion placing the reduced base inside the reduced image, cited from the same chapter and the algebra chapter. Cited at proof lines 11946-11949. Lets the argument for the claim be carried out after replacing the rings by their reductions, where the inclusion of the reduced base into the reduced image is available.
\item The structure result for a finite type algebra inducing a universal homeomorphism: generators satisfying square or cube relations, or p-th power relations, cited from the same chapter. Cited at proof lines 11954-11958. Supplies the finite list of generators and relations whose preservation at large indices is what makes the transition algebras induce universal homeomorphisms.
\item The algebra lemmas that maps of p-th power type and of 2-3 type induce universal homeomorphisms on spectra, cited from the algebra chapter. Cited at proof lines 11976-11978. Concludes that each step of the chain of intermediate algebras, and hence the whole map, induces a universal homeomorphism on spectra.
\end{itemize}\medskip
\ifdefined\hypertarget\hypertarget{sourcepointer-1}{}\fi
\noindent\textit{Source reference.} Original lemma 0EUH: focused-seeds/AGP/sources/stacks-snapshot-20261011/tags.txt, source lines 18552--18552. Same-chapter reference used in the retained author text at source line(s) 11974. It is supplied only as a navigation pointer so that the retained text compiles; the counted item delivers the author's own statement and the author's complete proof, and no part of the referenced item is copied in or substituted for the counted argument.\par\medskip
\ifdefined\hypertarget\hypertarget{sourcepointer-2}{}\fi
\noindent\textit{Source reference.} Original lemma 00QO: focused-seeds/AGP/sources/stacks-snapshot-20261011/tags.txt, source lines 974--974. Same-chapter reference used in the retained author text at source line(s) 10408, 11936. It is supplied only as a navigation pointer so that the retained text compiles; the counted item delivers the author's own statement and the author's complete proof, and no part of the referenced item is copied in or substituted for the counted argument.\par\medskip
\ifdefined\hypertarget\hypertarget{sourcepointer-3}{}\fi
\noindent\textit{Source reference.} Original lemma 00FL: focused-seeds/AGP/sources/stacks-snapshot-20261011/tags.txt, source lines 575--575. Same-chapter reference used in the retained author text at source line(s) 1301, 11842, 11946, 12203, 14779. It is supplied only as a navigation pointer so that the retained text compiles; the counted item delivers the author's own statement and the author's complete proof, and no part of the referenced item is copied in or substituted for the counted argument.\par\medskip
\ifdefined\hypertarget\hypertarget{sourcepointer-4}{}\fi
\noindent\textit{Source reference.} Original lemma 00I5: focused-seeds/AGP/sources/stacks-snapshot-20261011/tags.txt, source lines 667--667. Same-chapter reference used in the retained author text at source line(s) 11820, 11840, 11973. It is supplied only as a navigation pointer so that the retained text compiles; the counted item delivers the author's own statement and the author's complete proof, and no part of the referenced item is copied in or substituted for the counted argument.\par\medskip
\ifdefined\hypertarget\hypertarget{sourcepointer-5}{}\fi
\noindent\textit{Source reference.} Original lemma 00QN: focused-seeds/AGP/sources/stacks-snapshot-20261011/tags.txt, source lines 973--973. Same-chapter reference used in the retained author text at source line(s) 11932. It is supplied only as a navigation pointer so that the retained text compiles; the counted item delivers the author's own statement and the author's complete proof, and no part of the referenced item is copied in or substituted for the counted argument.\par\medskip
\ifdefined\hypertarget\hypertarget{sourcepointer-6}{}\fi
\noindent\textit{Source reference.} Original lemma 07C3: focused-seeds/AGP/sources/stacks-snapshot-20261011/tags.txt, source lines 9363--9363. Same-chapter reference used in the retained author text at source line(s) 11927. It is supplied only as a navigation pointer so that the retained text compiles; the counted item delivers the author's own statement and the author's complete proof, and no part of the referenced item is copied in or substituted for the counted argument.\par\medskip
\ifdefined\hypertarget\hypertarget{sourcepointer-7}{}\fi
\noindent\textit{Source reference.} Original lemma 054M: focused-seeds/AGP/sources/stacks-snapshot-20261011/tags.txt, source lines 6645--6645. Same-chapter reference used in the retained author text at source line(s) 11413, 11943, 14761. It is supplied only as a navigation pointer so that the retained text compiles; the counted item delivers the author's own statement and the author's complete proof, and no part of the referenced item is copied in or substituted for the counted argument.\par\medskip
\ifdefined\hypertarget\hypertarget{sourcepointer-8}{}\fi
\noindent\textit{Source reference.} Original lemma 0CN7: focused-seeds/AGP/sources/stacks-snapshot-20261011/tags.txt, source lines 15882--15882. Same-chapter reference used in the retained author text at source line(s) 11625, 11673, 11931. It is supplied only as a navigation pointer so that the retained text compiles; the counted item delivers the author's own statement and the author's complete proof, and no part of the referenced item is copied in or substituted for the counted argument.\par\medskip
\ifdefined\hypertarget\hypertarget{sourcepointer-9}{}\fi
\noindent\textit{Source reference.} Original proposition 0CNE: focused-seeds/AGP/sources/stacks-snapshot-20261011/tags.txt, source lines 15889--15889. Same-chapter reference used in the retained author text at source line(s) 11950, 14781, 14808, 14874. It is supplied only as a navigation pointer so that the retained text compiles; the counted item delivers the author's own statement and the author's complete proof, and no part of the referenced item is copied in or substituted for the counted argument.\par\medskip
\begin{lemma}
\label{lemma-universal-homeo-limit}
Let $A \to B$ be a ring map.
\begin{enumerate}
\item If $A \to B$ induces a universal homeomorphism on spectra,
then $B = \colim B_i$ is a filtered colimit of finitely presented $A$-algebras
$B_i$ such that $A \to B_i$ induces a universal homeomorphism on spectra.
\item If $A \to B$ induces isomorphisms on residue fields and
a universal homeomorphism on spectra, then $B = \colim B_i$ is a
filtered colimit of finitely presented $A$-algebras $B_i$ such that
$A \to B_i$ induces isomorphisms on residue fields and a
universal homeomorphism on spectra.
\end{enumerate}
\end{lemma}

\begin{proof}
Proof of (1). We will use the criterion of
Algebra, Lemma \ref{algebra-lemma-when-colimit}.
Let $A \to C$ be of finite presentation and let $\varphi : C \to B$
be an $A$-algebra map. Let $B' = \varphi(C) \subset B$ be the image.
Then $A \to B'$ induces a universal homeomorphism on spectra by
Lemma \ref{lemma-subalgebra-inherits}.
By Algebra, Lemma \ref{algebra-lemma-ring-colimit-fp}
we can write $B' = \colim_{i \in I} B_i$
with $A \to B_i$ of finite presentation and surjective
transition maps.
By Algebra, Lemma \ref{algebra-lemma-characterize-finite-presentation}
we can choose an index $0 \in I$
and a factorization $C \to B_0 \to B'$ of the map $C \to B'$.
We claim that $\Spec(B_i) \to \Spec(A)$ is a universal homeomorphism
for $i$ sufficiently large. The claim finishes the proof of (1).

\medskip\noindent
Proof of the claim. By Lemma \ref{lemma-reduction-universal-homeomorphism}
the ring map $A_{red} \to B'_{red}$ induces a universal homeomorphism
on spectra. Thus $A_{red} \subset B'_{red}$ by
Algebra, Lemma \ref{algebra-lemma-image-dense-generic-points}.
Setting $A' = \Im(A \to B')$ we have surjections $A \to A' \to A_{red}$
inducing bijections $\Spec(A_{red}) = \Spec(A') = \Spec(A)$.
Thus $A' \subset B'$ induces a universal homeomorphism on spectra.
By Proposition \ref{proposition-universal-homeomorphism}
and the fact that $B'$ is finite type over $A'$ we can
find $n$ and $b'_1, \ldots, b'_n \in B'$ such that
$B' = A'[b'_1, \ldots, b'_n]$
and such that for $j = 1, \ldots, n$ we have
\begin{enumerate}
\item $(b'_j)^2, (b'_j)^3 \in A'[b'_1, \ldots, b'_{j - 1}]$, or
\item there exists a prime number $p$ with
$pb'_j, (b'_j)^p \in A'[b'_1, \ldots, b'_{j - 1}]$.
\end{enumerate}
Choose $b_1, \ldots, b_n \in B_0$ lifting $b'_1, \ldots, b'_n$.
For $i \geq 0$ denote $b_{j, i}$ the image of $b_j$ in $B_i$.
For large enough $i$ we will have for $j = 1, \ldots, n$
\begin{enumerate}
\item $b_{j, i}^2, b_{j, i}^3 \in A_i[b_{1, i}, \ldots, b_{j - 1, i}]$, or
\item there exists a prime number $p$ with
$pb_{j, i}, b_{j, i}^p \in A_i[b_{1, i}, \ldots, b_{j - 1, i}]$.
\end{enumerate}
Here $A_i \subset B_i$ is the image of $A \to B_i$. Observe that $A \to A_i$
is a surjective ring map whose kernel is a locally nilpotent ideal.
After increasing $i$ more if necessary, we may assume $B_i$
is generated by $b_1, \ldots, b_n$ over $A_i$, in other words
$B_i = A_i[b_1, \ldots, b_n]$.
By Algebra, Lemmas \ref{algebra-lemma-p-ring-map} and
\ref{algebra-lemma-2-3-ring-map} we conclude that
$A \to A_i \to A_i[b_1] \to \ldots \to A_i[b_1, \ldots, b_n] = B_i$
induce universal homeomorphisms on spectra. This finishes the proof
of the claim.

\medskip\noindent
The proof of (2) is exactly the same.
\end{proof}

\end{document}
